\section{Números de Precisão Finita}
\title{Aula 5 - Precisão Finita}



\begin{frame}
    \titlepage
\end{frame}


\begin{frame}{Números de Precisão Finita}
    \begin{itemize}
        \item \textbf{No papel:} Usamos quantos dígitos forem necessários.
        \item \textbf{No Computador:} O espaço (memória) é limitado e definido no projeto do hardware.
    \end{itemize}

    \begin{block}{Definição}
        Computadores lidam com uma quantidade \textbf{fixa} de dígitos. Isso nos força a trabalhar com um conjunto limitado de valores representáveis.
    \end{block}

    \vspace{0.3cm}
    \textbf{O que estudaremos nesta seção:}
    \begin{itemize}
        \item Propriedades algébricas na precisão finita.
        \item Cálculo da quantidade de valores representáveis na base $B$.
        \item Representação de números inteiros com sinal.
    \end{itemize}
\end{frame}

\begin{frame}{Limitações da Precisão Finita}
    Considere um sistema hipotético que usa apenas \textbf{3 dígitos decimais} (000 a 999):

    \begin{itemize}
        \item \textbf{Valores impossíveis de representar:}
              \begin{itemize}
                  \item Maiores que 999 (Ex: 1000, 1500).
                  \item Menores que 0 (Ex: -1, -50).
                  \item Frações e números irracionais (Ex: 0,5 ou $\pi$).
              \end{itemize}
    \end{itemize}

    \begin{alertblock}{Quebra de Propriedades}
        Diferente da matemática tradicional, os números de precisão finita \textbf{não possuem a propriedade de fechamento}.
    \end{alertblock}
\end{frame}

\begin{frame}{A Propriedade de Fechamento}
    \begin{itemize}
        \item \textbf{Na Matemática (Conjunto $\mathbb{Z}$):} Para qualquer $i, j$, os resultados de $i+j, i-j$ e $i \times j$ também são inteiros.
        \item \textbf{No Computador:} O resultado de uma operação pode "escapar" da faixa permitida.
    \end{itemize}

    \begin{exampleblock}{Exemplos de Falha (Sistema de 3 dígitos)}
        \begin{itemize}
            \item \textbf{Soma:} $900 + 500 = 1400$ \hfill \textit{(Lupa: Cabe apenas 400?)}
            \item \textbf{Subtração:} $1 - 2 = -1$ \hfill \textit{(Não há sinal negativo)}
            \item \textbf{Multiplicação:} $100 \times 10 = 1000$ \hfill \textit{(Estourou o limite)}
        \end{itemize}
    \end{exampleblock}
\end{frame}

\begin{frame}{O Erro de Overflow}
    \begin{center}
        \Large \color{red} \textbf{OVERFLOW}
    \end{center}

    \begin{itemize}
        \item \textbf{O que é:} Ocorre quando o resultado de uma operação aritmética está fora da faixa de valores que o sistema consegue representar.
        \item \textbf{Analogia:} É como um hodômetro de carro que chega em 999.999km e volta para 000.000km.
    \end{itemize}

    \begin{block}{Impacto no Software}
        Se o programador não previr o overflow, o computador pode produzir resultados bizarros (como somar dois números positivos e obter um número negativo ou um valor muito pequeno).
    \end{block}
\end{frame}

\begin{frame}{Propriedades Algébricas na Precisão Finita}
    Diferente da Matemática Clássica, onde os conjuntos numéricos são infinitos, o hardware precisa lidar com limites físicos. Isso afeta como as propriedades algébricas se comportam.

    \begin{block}{Detecção de Erro}
        No caso de \textbf{overflow}, o hardware (CPU) é projetado para identificar o estouro de capacidade e emitir um aviso (\textit{flag}) para o sistema operacional.
    \end{block}

    \vspace{0.4cm}
    Vamos analisar as três propriedades fundamentais:
    \begin{enumerate}
        \item \textbf{Comutatividade}
        \item \textbf{Associatividade}
        \item \textbf{Distributividade}
    \end{enumerate}
\end{frame}

\begin{frame}{Comutatividade e Associatividade}
    \begin{itemize}
        \item \textbf{Comutatividade:} Mantém-se inalterada.
              \begin{itemize}
                  \item $a + b = b + a$
                  \item $a \times b = b \times a$
                  \item A ordem dos operandos não altera o resultado, mesmo que ocorra overflow.
              \end{itemize}

              \vspace{0.3cm}

        \item \textbf{Associatividade:} É satisfeita, \textbf{desde que o resultado final seja válido}.
              \begin{itemize}
                  \item $(a + b) + c = a + (b + c)$
                  \item $(a \times b) \times c = a \times (b \times c)$
              \end{itemize}
    \end{itemize}

    \begin{alertblock}{Associatividade}
        Se o resultado final estiver dentro da capacidade de representação, o valor será o mesmo. Se houver overflow, ele ocorrerá independentemente da ordem de execução, tornando o resultado inválido em ambos os casos.
    \end{alertblock}
\end{frame}


\begin{frame}{Distributividade na Precisão Finita}
    A última propriedade fundamental é a \textbf{Distributividade} da multiplicação sobre a adição:

    \begin{block}{Fórmula}
        $$ a \times (b + c) = (a \times b) + (a \times c) $$
    \end{block}

    \begin{itemize}
        \item \textbf{No Computador:} A propriedade se mantém válida para números inteiros.
        \item \textbf{A Ressalva:} Assim como na associatividade, a igualdade é garantida se:
              \begin{enumerate}
                  \item O resultado final for \textbf{válido} (dentro da faixa representável).
                  \item \textbf{OU} se ambos os caminhos de cálculo resultarem em um valor \textbf{inválido} (Overflow).
              \end{enumerate}
    \end{itemize}

    \begin{alertblock}{Conclusão das Propriedades}
        Para \textbf{Inteiros}, a Álgebra de Precisão Finita tenta mimetizar a Álgebra Clássica. O maior "vilão" não é a quebra das regras de ordem, mas sim o \textbf{estouro do limite de bits}.
    \end{alertblock}
\end{frame}

\begin{frame}{Por que as Propriedades importam?}
    Mesmo sem decorar os nomes, usamos essas regras para simplificar cálculos no dia a dia:

    \begin{itemize}
        \item \textbf{Comutatividade + Associatividade:} Permitem somar uma lista de números em qualquer ordem (ex: agrupar números que somam 10 para facilitar a conta de cabeça).
        \item \textbf{Distributividade:} É a base para \textbf{colocar fatores em evidência} em expressões algébricas.
    \end{itemize}

    \begin{block}{No Computador}
        Essas manipulações são usadas pelos \textbf{Compiladores} para otimizar o código, reorganizando as operações para que o processador as execute de forma mais rápida.
    \end{block}
\end{frame}

\subsection{Capacidade de Representação}

\begin{frame}{Quantos números cabem no sistema?}
    Para entender a capacidade de um hardware, precisamos resolver um problema de contagem:

    \begin{exampleblock}{O Problema}
        Dados $B$ símbolos diferentes, de quantas maneiras podemos preencher $n$ posições distintas (onde a ordem importa e símbolos podem se repetir)?
    \end{exampleblock}

    \begin{itemize}
        \item \textbf{Posição 0:} $B$ opções.
        \item \textbf{Posição 1:} $B$ opções.
        \item \dots
        \item \textbf{Posição $n-1$:} $B$ opções.
    \end{itemize}

    \begin{center}
        \Large \textbf{Total de Combinações = $B^n$}
    \end{center}
\end{frame}

\begin{frame}{Princípio Multiplicativo: $B^n$}
    A fórmula $B^n$ surge do princípio fundamental da contagem:

    \begin{itemize}
        \item Com \textbf{1 posição}: $B$ possibilidades.
        \item Com \textbf{2 posições}: $B \times B = B^2$ possibilidades.
        \item Com \textbf{3 posições}: $B \times B \times B = B^3$ possibilidades.
    \end{itemize}



    \begin{block}{Exemplo em Binário ($B=2$)}
        Se um computador usa **8 bits** ($n=8$):
        $$ 2^8 = 256 \text{ combinações diferentes} $$
        Se ele usa 10 bits:
        $$ 2^{10} = 1.024 \text{ combinações} $$
    \end{block}
\end{frame}

\begin{frame}{Exercícios: Capacidade de Representação}
    Utilizando a fórmula $B^n$, determine a quantidade de combinações possíveis para cada caso:


    \begin{enumerate}[a)]
        \item Quantos números diferentes podem ser expressos usando \textbf{2 dígitos} na base \textbf{10}?

              \vspace{0.4cm}

        \item Quantos números diferentes podem ser expressos usando \textbf{5 dígitos} na base \textbf{2}?

              \vspace{0.4cm}

        \item Quantos números diferentes podem ser expressos usando \textbf{32 dígitos} na base \textbf{2}?
    \end{enumerate}


    \vspace{0.6cm}
    \begin{center}
        \textit{Dica: Lembre-se da nossa tabela de potências de 2 vista anteriormente!}
    \end{center}
\end{frame}

\begin{frame}{Gabarito Comentado}
    \begin{itemize}
        \item[\textbf{a)}] \textbf{Base 10, $n=2$:}
              $$ 10^2 = \mathbf{100} \text{ combinações} $$
              \textit{(Os números de 00 a 99)}

              \vspace{0.3cm}

        \item[\textbf{b)}] \textbf{Base 2, $n=5$:}
              $$ 2^5 = \mathbf{32} \text{ combinações} $$
              \textit{(Os números binários de 00000 a 11111)}

              \vspace{0.3cm}

        \item[\textbf{c)}] \textbf{Base 2, $n=32$:}
              $$ 2^{32} = \mathbf{4.294.967.296} \text{ combinações} $$
              \textit{(Aproximadamente 4,3 bilhões de valores diferentes)}
    \end{itemize}


    A quantidade de bits ($n$) define o "tamanho do universo" de números que podemos contar. Se o resultado de uma conta for o número $2^{32} + 1$, ele não "caberá" no universo que so temos 32 bits, gerando o \textbf{Overflow}.

\end{frame}

\subsection{Convenções para Representação Finita}

\begin{frame}{Como representar o sinal?}
    Até agora, vimos números "sem sinal" (unsigned), onde todos os bits representam magnitude. Para incluir números negativos, precisamos de convenções:

    \begin{itemize}
        \item \textbf{Números sem sinal:} Faixa de $0$ a $B^n - 1$.
        \item \textbf{Números com sinal:} Três esquemas principais:
              \begin{enumerate}
                  \item \textbf{Sinal-Magnitude:} O bit mais à esquerda indica o sinal ($0 = +, 1 = -$).
                  \item \textbf{Excesso de $2^{n-1}$:} O valor é deslocado por uma constante.
                  \item \textbf{Complementos (1 e 2):} O bit de sinal faz parte do valor aritmético.
              \end{enumerate}
    \end{itemize}

    \begin{block}{O problema do Zero}
        Como $n$ bits geram uma quantidade par de combinações ($2^n$), ao usarmos uma para o zero, sobra um número ímpar de padrões para dividir entre positivos e negativos.
    \end{block}
\end{frame}

\subsection{Sinal-Magnitude}

\begin{frame}{O Esquema Sinal-Magnitude}
    É a forma mais intuitiva, semelhante ao que fazemos no papel:
    \vspace{0.3cm}
    \begin{center}
        \begin{tabular}{|c|c|c|c|c|}
            \hline
            $b_{n-1}$      & $b_{n-2}$                                                & \dots & $b_1$ & $b_0$ \\ \hline
            \textbf{Sinal} & \multicolumn{4}{c|}{\textbf{Magnitude (Valor Absoluto)}}                         \\ \hline
        \end{tabular}
    \end{center}

    \begin{itemize}
        \item \textbf{Bit de Sinal ($b_{n-1}$):} $0 \rightarrow$ Positivo ($+$) | $1 \rightarrow$ Negativo ($-$).
        \item \textbf{Magnitude:} Representada pelos $n-1$ bits restantes.
        \item \textbf{Faixa de Valores:} De $-(B^{n-1}-1)$ até $B^{n-1}-1$.
    \end{itemize}

    \begin{alertblock}{O Problema do "Zero Duplo"}
        Neste sistema, o zero possui duas representações:
        \begin{itemize}
            \item $+0$ (ex: \texttt{0000}) e $-0$ (ex: \texttt{1000}).
        \end{itemize}
        Isso desperdiça uma combinação e complica a lógica do hardware.
    \end{alertblock}
\end{frame}

\begin{frame}{Exemplos e Exercício (Sinal-Magnitude)}
    Considere $n=4$ bits:
    \begin{itemize}
        \item $0011_2 \rightarrow +3_{10}$
        \item $1011_2 \rightarrow -3_{10}$
        \item $1000_2 = 0000_2 = 0$
    \end{itemize}

    \vspace{0.4cm}

    Converta os decimais para binário usando \textbf{5 bits} em \textbf{Sinal-Magnitude}:
    \begin{enumerate}[a)]
        \item $-15$
        \item $-1$
        \item $0$ (apresente as duas formas)
        \item $10$
    \end{enumerate}

\end{frame}

\begin{frame}{Gabarito: Exercício 21}
    Para $n=5$ bits, o primeiro bit é o sinal e os outros 4 são a magnitude:

    \begin{itemize}
        \item[\textbf{a)}] $-15 \rightarrow$ Sinal: $1$ $\rightarrow$ Magnitude ($15_{10}$): $1111_2$ $\rightarrow$ \textbf{11111}
        \item[\textbf{b)}] $-1 \rightarrow$ Sinal: $1$ $\rightarrow$ Magnitude ($1_{10}$): $0001_2$ $\rightarrow$ \textbf{10001}
        \item[\textbf{c)}] $0 \rightarrow$ \textbf{00000} ($+0$) e \textbf{10000} ($-0$)
        \item[\textbf{d)}] $10 \rightarrow$ Sinal: $0$ $\rightarrow$ Magnitude ($10_{10}$): $1010_2$ $\rightarrow$ \textbf{01010}
    \end{itemize}

    \begin{block}{Reflexão}
        Note que em Sinal-Magnitude, para tornar um número negativo, basta "ligar" o bit mais significativo. O restante do número permanece idêntico.
    \end{block}
\end{frame}

\subsection{Excesso de $2^{n-1}$ (Representação com Viés)}

\begin{frame}{O Esquema de Excesso de $2^{n-1}$}
    Neste sistema, não "reservamos" um bit de sinal isolado. Em vez disso, somamos um valor fixo ao número real para deslocá-lo para a faixa positiva.

    \begin{block}{Funcionamento}
        Para $n$ bits, o excesso é $2^{n-1}$.
        $$\text{Representação Binária} = \text{Valor Decimal} + 2^{n-1}$$
    \end{block}

    \begin{itemize}
        \item \textbf{Inversão de Lógica:} Ao contrário do Sinal-Magnitude, aqui:
              \begin{itemize}
                  \item Números que começam com \textbf{0} são \textbf{Negativos}.
                  \item Números que começam com \textbf{1} são \textbf{Positivos} (ou zero).
              \end{itemize}
        \item \textbf{Vantagem:} Mantém a \textbf{ordem crescente} dos bits igual à ordem dos valores decimais.
        \item \textbf{Zero Único:} O valor 0 é representado pelo padrão $100...0_2$.
    \end{itemize}
\end{frame}

\begin{frame}{Exemplo: $n = 4$ (Excesso de $2^3 = 8$)}
    Com 4 bits, somamos 8 a cada valor para achar seu código:

    \begin{table}[]
        \centering
        \begin{tabular}{|c|c|l|}
            \hline
            \textbf{Decimal} & \textbf{Cálculo ($+8$)} & \textbf{Binário}       \\ \hline
            $-8$ (mínimo)    & $-8 + 8 = 0$            & \texttt{0000}          \\ \hline
            $-7$             & $-7 + 8 = 1$            & \texttt{0001}          \\ \hline
            \textbf{0}       & \textbf{$0 + 8 = 8$}    & \textbf{\texttt{1000}} \\ \hline
            $+1$             & $1 + 8 = 9$             & \texttt{1001}          \\ \hline
            $+7$ (máximo)    & $7 + 8 = 15$            & \texttt{1111}          \\ \hline
        \end{tabular}
    \end{table}

    \begin{itemize}
        \item \textbf{Faixa de Valores ($n=4$):} de $-8$ a $+7$.
        \item \textbf{Regra Geral:} de $-(2^{n-1})$ a $2^{n-1}-1$.
    \end{itemize}
\end{frame}



\begin{frame}{Exercício 22: Excesso de $2^{n-1}$ ($n=5$)}
    Para $n=5$, o excesso é $2^{5-1} = 2^4 = \mathbf{16}$.\\
    \textit{Regra: Some 16 ao número e converta o resultado para binário de 5 bits.}


    Converta para binário (5 bits, com Excesso de $2^{n-1}$):
    \begin{enumerate}[a)]
        \item $-15$
        \item $-1$
        \item $0$
        \item $10$
    \end{enumerate}
\end{frame}

\begin{frame}{Gabarito: Exercício 22}
    \begin{itemize}
        \item[\textbf{a)}] \textbf{$-15$:} $-15 + 16 = 1 \rightarrow \mathbf{00001_2}$
        \item[\textbf{b)}] \textbf{$-1$:} $-1 + 16 = 15 \rightarrow \mathbf{01111_2}$
        \item[\textbf{c)}] \textbf{$0$:} $0 + 16 = 16 \rightarrow \mathbf{10000_2}$
        \item[\textbf{d)}] \textbf{$10$:} $10 + 16 = 26 \rightarrow \mathbf{11010_2}$
    \end{itemize}

    \begin{exampleblock}{Dica Visual}
        Repare que o zero ($10000$) divide a tabela exatamente no meio. Qualquer número menor que $16$ começará com $0$ (negativos), e qualquer número $\geq 16$ começará com $1$ (positivos).
    \end{exampleblock}
\end{frame}